\chapter{细胞膜：一层会思考的边界吗}

\begin{kaichang}
想想洗手这件事。疫情那几年，你一定被叮嘱过：肥皂洗手，搓够 20 秒。奇怪的是，指令里常常不是“消毒液”，就是最普通的肥皂。为什么偏偏是肥皂？更奇怪的是，医生会告诉你：肥皂和 75\% 酒精对付某些病毒很管用，但对另一些病毒却效果平平——同样是病毒，待遇为什么不同？

再看一眼你的厨房：一块红甜菜根切开泡进冷水，水几乎不变色；泡进热水，几分钟就染成一片紫红。色素本来好好地关在细胞里，热水做了什么，把它们全放了出来？

第 25 章结尾我们留过一个悬念：鸡蛋膜是半透膜，水能过、溶质不能随心所欲地过，而活细胞的膜远比鸡蛋膜高明——上面装着“水泵”和“闸门”。这一章就兑现那个承诺。先猜猜看：一层只有\textbf{两个分子厚}的膜，凭什么既能挡住千军万马，又能精确放行特定分子，甚至还参与产生神经的电信号？它是一层会思考的边界吗？
\end{kaichang}

\section{现象层：看不见的边界，看得见的后果}

你从来没见过细胞膜。普通光学显微镜下，它只是细胞边缘一圈若有若无的轮廓——因为它只有大约 4 到 5 纳米厚，而光学显微镜的分辨极限大约是 200 纳米。一根头发丝的直径，够并排码放上万层细胞膜。

但这层看不见的边界，后果随处可见：

\begin{itemize}[itemsep=2pt]
\item 把红细胞泡进纯水，它会吸水胀破，鲜红的血红蛋白漏出来，悬液变得透明——这就是第 25 章算过的渗透压在动手：纯水一侧的水分子净涌入细胞，而细胞膜能承受的膨胀有限。
\item 把红细胞泡进浓盐水，它又失水皱缩，像一颗干瘪的葡萄干。只有泡在 0.9\% 的生理盐水里，它才安然无恙——等渗，不多不少。
\item 你吃进肚子的葡萄糖，几分钟到几小时内就出现在血液里，又被全身细胞飞快地取走。葡萄糖分子比水大得多，穿不过普通的脂质膜，细胞却能源源不断地把它运进门。
\item 你捏一下滚烫的杯子，手指几乎瞬间缩回。这个信号在神经里以每秒几十米的速度奔跑，而它的本质，是带电离子一波接一波地冲过神经细胞的膜。
\end{itemize}

挡水不挡，挡糖又放行，还能发电——这层的性格远比鸡蛋膜复杂。要理解它，得先把镜头推进到分子尺度，看看膜到底是用什么搭起来的。

\section{粒子层：会自动排队的磷脂}

膜的主要建筑材料，你在第 47 章已经认识过：磷脂。它长着一个亲水的“头”（带磷酸基团，爱和水打交道）和两条疏水的“尾”（脂肪酸长链，见水就躲）。当时我们说过，正是这种“两头性格分裂”的结构，让磷脂一遇水就自动组装成双层：尾巴对尾巴藏在内侧躲水，头朝外两面亲水——不需要图纸，不需要工头，纯属疏水效应推动的自组装。

现在请认真体会这件事有多不可思议。你把磷脂撒进水里，它们会\textbf{自己}围出一个个空心小球，把一小汪水包在内部，与外界隔开。换句话说：\textbf{边界，是这堆分子自己长出来的}。这件事一点也不平凡。第 45 章说过，生命系统的第一个特征就是“有边界”——有了边界，才有“里面”和“外面”，细胞才能维持一个和外界不一样的内部环境。而磷脂双层恰好提供了这种边界：它允许水和少数小分子穿过，却把大多数离子、糖、氨基酸挡在外面，或者至少让它们过得极其艰难。

这层膜有多薄？双层的总厚度约 4 到 5 纳米，就是两个磷脂分子首尾相接的长度。做个数量级估算，感受一下“薄”和“多”：一个红细胞大约是两面凹的圆盘，表面积约 140 平方微米，即 $1.4\times 10^{-10}$ 平方米；一个磷脂分子头部占的面积约为 $0.7$ 平方纳米，即 $0.7\times 10^{-18}$ 平方米。膜是双层、有内外两个表面，所以磷脂分子总数约为

\[2\times \frac{1.4\times 10^{-10}}{0.7\times 10^{-18}} = 4\times 10^{8}\]

约 \textbf{4 亿个磷脂分子}，才够糊住一个小小的红细胞。而你的全身大约有 20 多万亿个细胞。膜虽然薄到光学显微镜都无能为力，却是个数量惊人的大工程。

还有个细节值得记住：膜不是均质的磷脂片。里面嵌着大量蛋白质（有的膜里蛋白质占的质量比脂质还多），点缀着胆固醇（第 47 章提过，它像个“调稠剂”，让膜在不同温度下保持合适的流动性），外表面还挂着糖链——第 46 章说过，那些糖链是细胞的“身份标签”，供别的细胞识别。这是一堵由好几种材料混砌、还在不断流动的墙。

“不断流动”不是修辞，是功能必需。膜太硬，镶嵌的蛋白质没法移动、变形、工作；膜太稀，又留不住该留的东西。生物会主动调节膜的“稠度”：耐寒的鱼和越冬的植物，会增加膜脂里不饱和脂肪酸的比例——回想第 47 章，不饱和链带弯折、堆不紧密，就像往油里掺了防冻剂，让膜在低温下依然柔软。你的细胞则在膜里掺胆固醇干类似的事。发烧、冻伤之所以伤害细胞，部分原因正是膜的流动性被温度推离了正常区间。

\section{规则层：谁能过，怎么过}

现在进入核心问题：膜上的交通规则。物质穿过细胞膜，归根结底只有三条路，区分它们的标准只有两个——\textbf{需不需要蛋白质帮忙}、\textbf{耗不耗能量}。

\textbf{第一条路：简单扩散。}小分子直接从磷脂双层的缝隙里挤过去，不需要任何帮助。氧气、二氧化碳、水（大部分）走这条路。方向永远是从浓度高的一侧到浓度低的一侧——顺梯度下坡，和第 25 章讲的扩散、渗透是同一个统计规律：没有谁在指挥，只是高浓度一侧撞过去的分子更多。渗透其实就是水穿膜的简单扩散，第 25 章已经掰开揉碎讲过，这里不再重复。

\textbf{第二条路：易化扩散。}葡萄糖、各种离子想穿膜，但磷脂双层的疏水内部不欢迎它们——糖太亲水，离子带电。它们需要膜蛋白搭的“专用通道”：葡萄糖有葡萄糖转运蛋白，水有水通道蛋白，各种离子有各自的离子通道。注意：这条路虽然要“帮手”，但\textbf{不耗能量}，方向依然是顺梯度下坡。蛋白质只是降低了穿膜的门槛——用第 49 章的语言说，它降低的是“活化能”，不改变方向。方向由浓度梯度决定，谁也改不了。

\textbf{第三条路：主动运输。}这才是真正硬气的一条路：逆着梯度，把物质从低浓度一侧抽到高浓度一侧，上坡。上坡就要能量，能量来自第 50 章讲的 ATP 水解与运输过程的偶联。最著名的例子是钠钾泵：它每水解一份 ATP，就把 3 个钠离子泵出细胞、2 个钾离子泵进细胞——注意，这两个方向都是逆着浓度梯度的。

为什么要费这个劲？因为“梯度本身就是财富”。细胞外钠多、细胞内钾多，这个不对称的分布是钠钾泵长年累月花钱（ATP）维持出来的。而维持它的回报是：当细胞需要时，只要打开通道让离子顺坡冲下来，就能瞬间释放一股电流。这就像把水抽到高处的水塔里存着，用时一开龙头就有自来水——梯度是细胞储存的“势能存款”。

三条路都在你身边日夜运转。你肾脏里的肾小管细胞靠水通道蛋白快速回收水分，一天过滤出的原尿约 180 升，被重吸收后只剩一两升尿；某些利尿药正是通过干扰离子在肾小管的重吸收，让水“留不住”地排出去。而通道一旦被堵住，就是中毒：河豚毒素能恰好卡进神经细胞钠通道的孔口，离子洪流断流，神经信号瘫痪——所以河豚料理必须持证上岗，这不是玄学，是膜蛋白几何形状的生死问题。

这三条路也不是各自孤立的。想象小肠里的一场接力：肠腔里葡萄糖浓度低、细胞里反而高，直接扩散是上坡，怎么办？细胞先把钠离子用钠钾泵逆坡泵出（耗 ATP），再让一种“共运输蛋白”放钠离子顺坡回流——同时硬性规定：钠必须拉着葡萄糖一起过闸。钠下坡释放的“势能”，就这样转手垫付了葡萄糖的上坡。葡萄糖不耗 ATP 却完成了主动运输，账其实是钠钾泵结的。你看，膜上的每一台机器都在遵循同一条总账：能量不凭空来，只能从一个梯度倒手到另一个梯度。

\begin{qiaoliang}
这笔存款有多厚？神经和肌肉细胞内外，钠离子浓度大约相差 10 倍，钾离子相差约 30 倍，再加上膜对钾的通透性远高于钠，净效果是膜两侧维持着一个约 $-70$ 毫伏的电压差（内负外正），叫\textbf{膜电位}。$-70$ 毫伏听起来微不足道，但别忘了膜只有约 5 纳米厚——把电压除以厚度，膜内的电场强度高达每米上千万伏，是高压输电线路周围电场的好几倍。只是因为距离太短，总能量很小，才不至于把你电着。一旦某个瞬间钠通道打开，钠离子顺坡涌入，局部电位在千分之一秒内从 $-70$ 毫伏翻到 $+30$ 毫伏——这就是神经信号的“电脉冲”，叫动作电位。你缩手的速度，取决于这波“离子洪水”沿神经传导的速度。
\end{qiaoliang}

\section{符号层：画一张膜的地图}

让我们把上面的一切浓缩成一张图。习惯上，生物学家用 $[\chem{Na^+}]_{\text{外}}$ 这样的记号表示“细胞外的钠离子浓度”，用箭头表示运输方向，用 $-70$ mV 表示膜电位。下面这幅示意图把三种运输方式画在了一起：

\begin{center}
\begin{tikzpicture}[scale=0.9]
\foreach \x in {0,1,2,3,4,5,6,7,8} {
  \fill[blue!60] (\x,0.45) circle (0.13);
  \draw[blue!60] (\x,0.35) -- (\x-0.07,-0.05);
  \draw[blue!60] (\x,0.35) -- (\x+0.07,-0.05);
  \fill[blue!60] (\x,-0.55) circle (0.13);
  \draw[blue!60] (\x,-0.45) -- (\x-0.07,-0.85);
  \draw[blue!60] (\x,-0.45) -- (\x+0.07,-0.85);
}
\draw[orange,fill=orange!30] (2.5,0.7) rectangle (3.5,-1.0);
\node at (3,-0.15) {\small 通道};
\draw[red,fill=red!20] (5.5,0.7) circle (0.35);
\draw[red,fill=red!20] (5.5,-1.0) circle (0.35);
\draw[red,thick] (5.85,0.7) -- (5.85,-1.0);
\node at (5.5,-1.6) {\small 泵};
\draw[->,thick] (3,1.3) node[above]{\small $\chem{Na^+}$ 顺坡} -- (3,0.75);
\draw[->,thick] (5.5,-1.45) -- (5.5,-2.1) node[below]{\small 逆坡泵出，耗 ATP};
\draw[->,thick,dashed] (0.8,1.3) node[above]{\small $\chem{O_2}$ 直接穿过} -- (0.8,-1.3);
\node at (-1.2,1.1) {\small 细胞外};
\node at (-1.2,-1.9) {\small 细胞内};
\end{tikzpicture}
\end{center}

照例声明：这是人为的表示法。圆球和短线不是磷脂的“照片”，真实的分子没有颜色，也在不断晃动、交换位置；蛋白质和脂质的相对大小也被严重压缩了——很多膜蛋白比双层本身还“高”，像冰山一样两端露出水面。

\section{证据层：没人见过的东西，怎么知道它长这样}

\begin{zhengju}
细胞膜薄到任何显微镜在半个多世纪里都看不见它，那科学家凭什么断定它是“磷脂双层”？答案是一串环环相扣、还翻过车的推理。

第一步在 1925 年。两位荷兰科学家——高特和格伦德尔——设计了一个聪明的实验。他们从红细胞中提取出全部脂质（选红细胞，是因为它内部没有细胞器，脂质几乎全来自那层外膜，账本干净）。然后把脂质滴在水面上，让它们铺展成一层单分子膜，测量这层膜的总面积。结果：脂质铺开的面积，恰好是红细胞表面积的\textbf{两倍}。结论几乎是唯一解：细胞膜由\textbf{两层}脂质分子构成。一个看不见的结构，被“数面积”数了出来。有趣的是，后来人们发现他们的实验其实有两个相互抵消的误差（脂质提取不完全，表面积也算小了），结论却歪打正着地对了——科学史并不总是伟光正的直线。

第二步却走了弯路。1935 年，戴维森和达尼埃利提出：双层两侧应该各盖着一层蛋白质，像三明治。这个“三明治模型”统治了教科书三十多年。但裂缝渐渐出现：电子显微镜下的膜确实是三层深色结构，看似支持三明治；可是按三明治模型，蛋白质把脂质完全盖住，那膜表面的化学性质应该处处一样，实际测量却不是；更要命的是，很多膜蛋白长着大段疏水的部分，根本不可能舒服地摊在膜表面。

1972 年，辛格和尼科尔森提出了\textbf{流动镶嵌模型}：蛋白质不是摊在表面的“盖子”，而是像冰山一样\textbf{镶嵌}在脂质海洋里，有的贯穿整个双层，有的只嵌入一半；脂质和蛋白质都能在膜平面内侧向漂移。关键证据之一来自“冰冻蚀刻”技术：把膜冻硬后劈开，断面常常沿双层正中间裂开，露出的断面上嵌着密密麻麻的颗粒——那些就是贯穿膜的蛋白质的内部断面。三明治模型预言断面应该光溜溜的，事实宣判了它的死刑。

故事还没完。后来的研究发现，膜并不是均匀的“海洋”：胆固醇和某些脂质会聚拢成一个个相对稳定的“小岛”（被称为脂筏），特定蛋白质爱在这些岛上扎堆。膜也不是自由液面——细胞内部的骨架纤维会“拴住”很多膜蛋白，让它们不能随意漂。流动镶嵌模型至今仍是最好的框架，但它已经和 1972 年的版本不太一样了。这就是证据故事的标准节奏：数出来的双层、推翻的三明治、立起来的流动镶嵌、打补丁的脂筏——每一步都有人提问题、做实验、找替代解释。
\end{zhengju}

\section{选择性响应：没有意识，却有“判断力”}

现在可以正面回答标题的问题了。细胞膜表现得极其“聪明”：葡萄糖来了，葡萄糖转运蛋白放行，结构相似的别的糖却进不去；神经递质飘来，特定受体被激活，细胞立刻做出反应；激素到了，细胞认得它，旁观的分子再多也视若无睹。这一切看起来像一个有判断力的门卫。

但请记住这本书的口头禅：别问粒子“想”做什么，只问哪条路在物理上行得通。膜的“判断力”全部来自三件毫无意识可言的东西：

\begin{itemize}[itemsep=2pt]
\item \textbf{形状与电荷的匹配。}转运蛋白和受体的结合部位有特定的几何形状和电荷分布，只有互补的分子能卡住——和第 48 章讲的蛋白质结构决定功能、第 49 章讲的诱导契合，是同一条原理。这不是“认出”，是“卡得上”。
\item \textbf{能量偶联。}主动运输逆坡上行的每一步，都由 ATP 水解（或现成的离子梯度）“垫付”能量。没有能量输入，泵立刻停摆——它没有任何“坚持工作”的意志。
\item \textbf{统计规律。}扩散和渗透的净方向，纯粹是两边碰撞次数的差。膜两侧没有任何信息传递，水分子不知道对面有多浓。
\end{itemize}

受体尤其能说明问题。你鼻子里有几百种嗅觉受体，每种受体的结合口袋形状不同，能“卡住”不同气味分子；卡住的那一刻，受体变形，触发细胞内一连串分子事件，最终打开离子通道、改变膜电位——你闻到了咖啡香。整个过程没有一步涉及“意识”，全是形状、能量和概率。所以标题的答案是：细胞膜\textbf{不会思考}。它的“选择”是分子结构与物理规律的必然结果，就像锁不“认识”钥匙，只是钥匙的形状恰好匹配。把一整套这样的机制叠加起来，才涌现出了看似智能的行为——这个“简单规则叠加出复杂表现”的主题，会在本书后面反复回响。

\begin{wuqu}
“细胞膜像智能保安，认识什么该进、什么不该进。”——反例俯拾皆是。一氧化碳分子形状和氧气有些相似，血红蛋白就“上当”把它抓了进来，这是煤气中毒的分子根源（第 15 章）；毒素乌本苷的形状恰好能卡住钠钾泵的某个位置，泵被活活卡死，细胞因梯度崩溃而死；铅离子能混过某些钙通道，因为电荷相同、大小相近。真要有智能保安，这些致命错误一个都不会发生。膜的“识别”只是形状和电荷的匹配游戏——匹配出错的代价，就是毒理学的全部内容。也正因如此，药物设计的核心思路之一，就是造出形状“差一点”的分子去骗过、堵住或劫持这些没有判断力的门卫。
\end{wuqu}

\section{边界层：这个模型哪里会漏}

流动镶嵌模型好用，但用之前先说好它的边界：

\begin{itemize}[itemsep=2pt]
\item “膜是均匀二维液体”只是近似。脂筏、细胞骨架的锚定、膜蛋白之间的聚集，都让真实的膜更像一座有街区、有拴系桩的城市，而不是一锅均匀的汤。研究膜蛋白在膜上的真实运动，至今是活跃的前沿。
\item “通道永远畅通”是错的。多数离子通道有“门”，受电压、配体或机械力的控制而开闭——神经信号的奥秘一半在通道，另一半就在这些门的开闭规律上。把通道想成一根固定的管子，会漏掉最精彩的部分。
\item 人工磷脂双层和真实细胞膜相去甚远。真实膜有成百上千种脂质、两面不对称（内外层的脂质成分不同，而且细胞会花能量维持这种不对称）、蛋白质密集到几乎“摩肩接踵”。实验室里干净的模型膜是简化，不是写真。
\item “渗透压公式”“梯度决定方向”这些定量结论，来自第 25 章就声明过的稀溶液近似。细胞质是拥挤的浓溶液，一切都只有定性意义上的“差不多”。
\end{itemize}

\begin{shiyan}
\textbf{热水怎样“拆掉”细胞的围墙}

材料：一小块红甜菜根（或紫甘蓝叶）、两个透明杯子、冷水、热水（约 $60\,^{\circ}\mathrm{C}$，饮水机热水档即可）、水果刀（请家长帮忙切）。

步骤：把甜菜根切成大小相近的两小块，自来水下轻轻冲掉切面上渗出的色素。分别投入冷水杯和热水杯，静置 10 分钟，期间不要搅动。

观察点：哪个杯子的水先变色？颜色深浅差多少？对比一下，热水里发生了什么？

原理：甜菜根的红色素平时被关在细胞的液泡里，液泡膜和细胞膜两道脂质双层把它锁得严严实实。加热会让磷脂双层过度流动、让膜蛋白变性（第 48 章讲过热变性），围墙出现窟窿，色素倾巢而出——你看到的颜色，就是膜的“阵亡通知书”。

安全提示：热水防烫，杯子选耐热的；实验材料接触过生水和刀具，做完不要食用。
\end{shiyan}

\section{回到开场：肥皂拆墙记}

现在可以回答洗手的问题了。包括新冠病毒、流感病毒、艾滋病毒在内的许多病毒，外面套着一层\textbf{包膜}——那其实就是一层磷脂双层，是病毒从宿主细胞膜上“偷”来的外套，上面插着病毒自己的蛋白质。这层外套对病毒至关重要：它既是保护壳，又是入侵下一个细胞时的“ docking 装置”（包膜先与细胞膜融合，病毒内容物才得以注入）。

而肥皂分子是什么？它和磷脂是亲戚：一个亲水头、一条疏水尾，只是只有一条尾巴。当肥皂分子涌向病毒包膜，它们的疏水尾巴会插进脂双层，亲水头露在外面——结果就是双层的自组装平衡被搅乱：原本严丝合缝的围墙被拆成一块块混着肥皂的碎片，病毒的内容物漏出来，病毒就此报废。你的甜菜根实验里热水干的事，肥皂在病毒身上干得更精准。75\% 酒精的思路类似：它能溶解脂质、同时让蛋白质变性，双管齐下。为什么不是纯酒精？纯酒精让蛋白质表面瞬间凝固，反而像给病毒“焊了层壳”，保护内部；有水在场的 75\% 配方渗透和变性才配合得最好——剂量与配比，又是化学。

但注意大纲里那个“不是对所有病毒都同样有效”。诺如病毒、脊髓灰质炎病毒、腺病毒这些\textbf{无包膜病毒}，外壳直接由蛋白质拼装而成，没有脂质双层可拆——肥皂对它们的杀伤力大打折扣（洗手时主要靠水流物理冲走），酒精也常常力不从心，对付它们更可靠的是含氯消毒剂（第 17 章讲过的那位暴躁的氯）。“酒精、肥皂能杀死一切病毒”的说法，错就错在没问病毒穿什么衣服。这再一次说明：了解结构，才能预测手段是否奏效——这正是全书五个贯穿问题里“它由什么组成、怎样连接”的实战。

\section{继续深挖：把药装进脂质体，以及下一章的胃口}

膜的原理早已不只是解释生命，还在被工程师反向使用。既然磷脂遇水自动围成小球，那就把药物装进小球里——这叫\textbf{脂质体}。脂质体的膜和细胞膜“材质相同”，能与细胞膜融合，把包裹的药物直接送进细胞内部。一些抗癌药就是这么投递的。而更出名的应用你多半接种过：mRNA 疫苗里，那段 mRNA 正是裹在脂质纳米颗粒里送进细胞的——没有这层人造膜，脆弱的 mRNA 连细胞的大门都摸不到。

医学上另一台“膜机器”是血液透析：肾功能衰竭的病人，血液流过半透膜做的透析器，代谢废物顺浓度梯度扩散出去，有用的物质被膜留下——一台体外的人工肾脏，用的全是本章的扩散与渗透原理，只是“选择透过”由膜的孔径机械决定，远不如活细胞的膜那样精细。

回头看这条因果链：第 47 章的磷脂结构 → 自组装成双层 → 双层上的运输规则 → 梯度与膜电位 → 神经信号、病毒攻防、药物递送。一块肥皂、一根神经、一剂疫苗，被同一条链串了起来。

而故事刚刚铺开。膜的存在，意味着细胞的“里面”是一个被精心维护的特殊环境：糖运进来了，离子各就各位了，ATP 的储备就位了。接下来的问题是：葡萄糖进了细胞大门之后，细胞怎么把它一点点拆开、把能量取出来？为什么不一把火烧掉，而要分成十几步、像拆炸弹一样小心翼翼？那是第 52 章《糖酵解》的故事。你还会在第 53 章看到膜再次登场——线粒体内膜上的质子梯度，是细胞发电站的真正心脏，那里的膜玩法比本章还要精妙。

\begin{tiaozhan}
膜电位可以用一条能斯特方程定量描述。某种离子在膜两侧达到电化学平衡时，平衡电位为 $E = \frac{RT}{zF}\ln\frac{c_{\text{外}}}{c_{\text{内}}}$，其中 $z$ 是离子电荷数，$F$ 是法拉第常数。代入体温、钾离子 $c_{\text{内}}\approx 150$ mM、$c_{\text{外}}\approx 5$ mM，算出钾的平衡电位约 $-90$ 毫伏——和实际膜电位 $-70$ 毫伏接近但不完全相等，差值来自膜对钠等离子也有少量通透性（需要用考虑多种离子的戈德曼方程修正）。真实的膜电位，是多种离子各自“想达到的电压”按通透性加权拉锯出来的折中结果。跳过本节不影响主线。
\end{tiaozhan}

\section*{练一练}

\begin{sikaoti}
\item 画一幅粒子示意图：用圆圈和短线画出一小段磷脂双层，在膜上画出“简单扩散、易化扩散、主动运输”三条路径，并各举一个走这条路的物质。图旁注明哪些步骤消耗 ATP，哪些只是顺坡下滑。
\item 把一株植物细胞分别放进纯水和浓盐水（植物有细胞壁，不会胀破）。用第 25 章的渗透原理预测两种情况下细胞会发生什么，并解释为什么腌菜用的浓盐水能让菜蔫掉。
\item 钠钾泵每循环一次：3 个钠离子出、2 个钾离子进，消耗 1 个 ATP。从电荷角度算一算：每循环一次，膜两侧净移出了几个正电荷？这对“细胞内偏负”的膜电位有什么贡献？
\item 有一种新发现的病毒，用肥皂水处理后完全失去感染能力，用蛋白酶处理后也失去感染能力。仅凭这两条证据，你能推断它的外壳结构吗？还需要排除哪些替代解释？
\item 为什么“纯酒精消毒效果反而不如 75\% 酒精”？请从膜和蛋白质两个角度分别解释。这个例子对你理解“剂量与配比”有什么启发？
\item 本章说“膜的识别只是形状匹配，不是认识”。请用这个观点解释：为什么误采毒蘑菇会中毒——毒素是“故意”攻击人体的吗？把答案写成两三句因果链。
\end{sikaoti}
